\documentclass[UTF8,fontset=none,11pt,a4paper]{ctexart}
\usepackage{ipara-handbook}
\handbooksetup
  {DS-01 线性关系与存储表示}
  {408 · 数据结构 DS-01}
  {Relation · Representation · Invariant · Cost}
\tikzset{
  cell/.style={draw=iparaDeep,minimum width=1.25cm,minimum height=0.72cm,align=center,fill=iparaSoft,font=\small},
  node/.style={draw=iparaGreen,rounded corners=1.2mm,minimum width=1.7cm,minimum height=0.72cm,align=center,fill=white,font=\small},
  flow/.style={-{Latex[length=2mm]},thick,draw=iparaDeep},
  ref/.style={-{Latex[length=2mm]},dashed,draw=iparaRule}
}
\begin{document}
\handbooktitle
  {线性关系与存储表示}
  {同一逻辑序列怎样换表示，操作为什么随之改变}
  {408 · 数据结构 Topic DS-01}

\begin{corebox}{母问题：逻辑线性关系怎样被有限存储兑现？}
线性表不是“数组或链表的名字”，而是一组对元素顺序和操作结果的契约。手册的生成链是：
\[
\begin{aligned}
\text{Linear Relation}&\to\text{Operation Contract}\to\text{Representation}\\
&\to\text{Local Update}\to\text{Invariant}\to\text{Cost Vector}
\end{aligned}
\]
箭头先表示建模顺序，再表示一次操作中信息和状态如何流动：先固定逻辑对象，才知道某种物理布局是否值得维护。不同表示必须给出相同的抽象结果，但不必承担相同的移动、寻址和额外空间成本。
\end{corebox}

\begin{methodbox}{本册定位与 Owner 边界}
\begin{itemize}
\item \kw{Owns：}线性表 ADT、顺序表、单链表、双链表、静态链式表示；数组的多维地址映射与特殊矩阵压缩；插入/删除/按位访问/按值查找的状态变化与边界。
\item \kw{Uses：}Atlas Foundation 的问题规模、基本操作、最坏/平均/摊还成本和成本向量语言。
\item \kw{Outputs：}向 DS02 提供“底层顺序/链接表示”，向 DS09 提供“有序序列上的访问与插入代价”；栈、队列、字符串和索引只调用本册表示，不复制本册机制。
\item \kw{Stops：}受限端点语义归 DS02；KMP 归 DS03；树/图节点结构归 DS04/DS07；BST、AVL、B/B+ 的有序索引机制归 DS09。
\end{itemize}
\end{methodbox}

\tableofcontents\newpage

\section{先把“表”与“存储”分开}
抽象线性表 $L=(a_1,a_2,\ldots,a_n)$ 只承诺元素的线性次序、长度和允许操作。它不承诺元素在内存中连续，也不承诺一次操作一定只访问一个地址。若把逻辑位置记为 $i$，必须先区分：
\begin{center}\small
\begin{tabularx}{\linewidth}{L{2.1cm}YY}
\toprule
层次 & 要回答的问题 & 不能直接推出的结论\\
\midrule
逻辑关系 & 谁是第 $i$ 个元素？插入后次序怎样？ & 元素一定相邻\\
操作契约 & 按位访问、按值查找、插入、删除的输入是什么？ & 每种操作都同样便宜\\
物理表示 & 数据、长度、容量、指针和空闲单元怎样保存？ & 只看元素字段就能判断状态\\
成本模型 & 计数对象是比较、移动、指针跳转还是 I/O？ & 一个 $O(\cdot)$ 覆盖所有场景\\
\bottomrule
\end{tabularx}\end{center}

\begin{examplebox}{最小反例：同一个“插入第 $k$ 个”}
对逻辑表 $[a,b,c,d]$ 在位置 $3$ 插入 $x$，抽象结果只有一个：$[a,b,x,c,d]$。顺序表需要把 $c,d$ 向后移动；链表只需让 $x$ 接入正确的指针链，但前提是已经拿到插入点的前驱。若题目只给关键字而不给位置，寻找前驱的成本又回来了。
\end{examplebox}

\section{顺序表：地址映射换取按位访问}
顺序表把元素放入一段可按下标解释的连续区域。设首地址为 $B$，每个元素占 $w$ 个存储单元，则零基下标 $i$ 的地址为
\[
\operatorname{addr}(A[i])=B+i\,w.
\]
这条式子不是“数组 API 口诀”，而是顺序表能在 $\Theta(1)$ 时间定位第 $i$ 个元素的原因：位置到地址是一次算术映射，不需要沿中间元素搜索。

\subsection{长度、容量与合法状态}
动态顺序表至少维护 $(data,\;length,\;capacity)$。合法状态包括
\[
0\le length\le capacity,\qquad 0\le i<length\Rightarrow A[i]\text{ 是有效元素}.
\]
插入前先检查 $length<capacity$；若需要扩容，旧元素必须按原次序复制到新区域，再改变容量和写入位置。扩容后的单次复制可能是 $\Theta(n)$，但若容量按固定倍数增长，连续追加的总复制量可摊还到每次 $O(1)$；这里的“摊还”不是把一次最坏成本改写成 $O(1)$。

\subsection{插入与删除是区域移动}
在长度为 $n$ 的表中，在逻辑位置 $k$（$0\le k\le n$）插入：
\[
\text{检查容量}\to A[n-1]\rightsquigarrow A[k]\text{ 右移}\to A[k]\leftarrow x\to length{+}{+}.
\]
需要移动 $n-k$ 个元素；在头部插入为 $\Theta(n)$，尾部追加在有剩余容量时为 $\Theta(1)$。删除位置 $k$ 则把区间 $A[k+1\ldots n-1]$ 左移，移动 $n-k-1$ 个元素，再递减长度。覆盖或清理被删除槽位是实现选择，但不能让无效槽位继续被逻辑遍历访问。

\begin{center}
\begin{tikzpicture}[node distance=0.25cm]
\node[cell](a0){a};\node[cell,right=of a0](a1){b};\node[cell,right=of a1](a2){c};\node[cell,right=of a2](a3){d};
\node[cell,right=2.2cm of a3](b0){a};\node[cell,right=of b0](b1){b};\node[cell,right=of b1](b2){x};\node[cell,right=of b2](b3){c};\node[cell,right=of b3](b4){d};
\node[above=0.18cm of a2,align=center]{插入前};\node[above=0.18cm of b2,align=center]{插入后};
\draw[flow] (a3.east)--(b0.west);
\draw[ref] (a2.south)--(b3.north);\draw[ref] (a3.south)--(b4.north);
\end{tikzpicture}
\end{center}

\section{多维数组与特殊矩阵：坐标如何映射为线性偏移}
数组的维数是逻辑坐标的个数，内存地址仍是一维序列。行优先、列优先不是两种数组对象，而是两种线性化函数。设二维数组 $A[L_1..U_1][L_2..U_2]$，每个元素占 $w$ 个存储单元，$N_1=U_1-L_1+1$，$N_2=U_2-L_2+1$，则
\[
\begin{aligned}
\operatorname{addr}_{row}(A[i][j])&=B+\bigl((i-L_1)N_2+(j-L_2)\bigr)w,\\
\operatorname{addr}_{col}(A[i][j])&=B+\bigl((j-L_2)N_1+(i-L_1)\bigr)w.
\end{aligned}
\]
两式都从同一问题生成：目标坐标之前，线性序列中已经排了多少个元素？下界不从 $0$ 或 $1$ 开始时，必须先减去逻辑下界；只背 $i n+j$ 会在非零下界题中错位。

对 $d$ 维数组，行优先偏移可递推为
\[
\operatorname{offset}=(((i_1-L_1)N_2+(i_2-L_2))N_3+\cdots)N_d+(i_d-L_d),
\]
本质是混合进位制，不需要另背展开式。

\subsection{压缩存储：先找独立元素，再建立可逆编号}
特殊矩阵只有在未存部分可由结构规则恢复时才能压缩。先确定独立元素集合，再为它建立连续编号 $k=f(i,j)$：
\begin{itemize}
\item 对称矩阵只存一个三角区，因为 $a_{ij}=a_{ji}$；访问另一半时先交换坐标再编号。
\item 上/下三角矩阵存一个三角区和一个公共常数；另一三角区的元素不占独立槽位。
\item 对角、带状矩阵只存满足带宽约束的坐标；带外元素由题设常数恢复。
\item 稀疏矩阵通常保存非零三元组 $(row,column,value)$；它换掉了连续随机定位能力，查找需额外搜索或索引。
\end{itemize}
例如按行存储下三角矩阵（含主对角线），零基坐标满足 $0\le j\le i<n$。第 $i$ 行之前已有 $i(i+1)/2$ 个元素，所以
\[
k=\frac{i(i+1)}2+j.
\]
公式证据是“前面完整行的长度和 + 当前行内偏移”。若改用一基下标或上三角按列存储，应重新数前置元素，不能机械平移。

\begin{methodbox}{数组映射题的停止条件}
先写合法坐标域、存储顺序和独立元素集合；再数目标元素之前的槽位，最后乘元素宽度并加首地址。编号必须落在 $[0,m-1]$，其中 $m$ 是实际槽数；若两个本应独立的坐标映到同一编号，映射就不合法。
\end{methodbox}

\section{链式表示：指针关系换取局部修改}
单链表节点保存 $(value,next)$，表对象通常保存头指针和可选的尾指针、长度。逻辑相邻不等于物理相邻：$a_i$ 的后继由指针给出，而不是由地址加 $w$ 得到。双链表再保存 $prev$，换取从后继回到前驱的能力。

\subsection{链表的不变量}
对无环单链表，必须保持：
\begin{enumerate}
\item 头指针要么为空，要么指向第一个有效节点；
\item 从头沿 $next$ 访问，每个有效节点至多出现一次，最终到达空指针；
\item 若维护 $tail$，则 $tail$ 是最后节点且 $tail.next=null$；
\item 若维护 $length$，它等于从头可达节点数。
\end{enumerate}
双链表还要保持相邻互逆：$u.next=v\Rightarrow v.prev=u$。任一方向只改一条边都会暂时破坏结构，后续操作可能沿错误方向丢失节点。

\subsection{局部插入的完整生命周期}
已知前驱 $p$，在其后插入新节点 $x$ 的最小动作是：
\[
x.next\leftarrow p.next,\qquad p.next\leftarrow x.
\]
若 $p$ 原来是尾节点，还要更新 $tail$；若向空表插入，还要同时更新 $head$。删除节点 $x$ 时，必须先拿到前驱 $p$，令 $p.next\leftarrow x.next$，再处理尾指针、长度和资源回收。\kw{“指针修改 $O(1)$”只描述已定位节点后的局部动作，不包括按关键字寻找节点。}

\section{同一操作的两套成本向量}
不要用一个复杂度数字抹平表示差异。对典型操作，至少记录“是否已知位置/前驱、比较数、移动或跳转数、额外空间和最坏边界”：
\begin{center}\small
\begin{tabularx}{\linewidth}{L{2.5cm}C{1.8cm}C{1.8cm}YY}
\toprule
操作 & 顺序表 & 单链表 & 成本解释 & 关键前提\\
\midrule
按位访问 $i$ & $\Theta(1)$ & $\Theta(n)$ & 地址映射 vs 指针跳转 & $i$ 合法\\
按值查找 & $\Theta(n)$ & $\Theta(n)$ & 逐项比较 & 无额外索引\\
已知位置插入 & $\Theta(n)$ & $\Theta(1)$ & 区间移动 vs 改两条边 & 链表已知前驱\\
头部插入/删除 & $\Theta(n)$ / $\Theta(n)$ & $\Theta(1)$ / $\Theta(1)$ & 顺序表整体移动；链表改头指针 & 正确维护 head\\
尾部追加 & $\Theta(1)$ 摊还 & $\Theta(1)$ & 容量复制 vs tail 指针 & 无扩容/已维护 tail\\
\bottomrule
\end{tabularx}\end{center}
空间也要分开：顺序表可能有未使用容量，链表每个节点额外保存指针；两者的逻辑元素空间都是 $\Theta(n)$，但常数、局部性和内存管理成本不同。

\section{生成性例子：把“按值删除”拆成两个阶段}
给定表 $L=[7,2,9,2,5]$，删除第一次出现的 $2$。这不是一个动作，而是：
\begin{enumerate}
\item \kw{Locate：}从头按逻辑次序比较，找到位置 $k=1$；
\item \kw{Update：}顺序表把 $A[2\ldots4]$ 左移，链表把前驱的后继越过目标节点；
\item \kw{Repair：}更新长度、尾指针、空槽或回收节点；
\item \kw{Verify：}再次遍历，确认次序为 $[7,9,2,5]$ 且所有结构不变量成立。
\end{enumerate}
两种表示的 Locate 都可能是 $\Theta(n)$；因此不能看到链表的局部删除就断言“链表删除总是 $O(1)$”。真正改变的是 Update 的成本和维护方式。

\section{代码把模型变成可执行状态转移}
代码不是本册结论之后的语法附录。它必须逐项兑现操作契约、状态字段、不变量修复和边界分支。完整 C++17 实现位于 \codepath{code/ds01_linear_list.hpp}；下面的代码块直接从该文件抽取，因此正文与可运行实现共享同一来源。

\subsection{顺序表：先保证容量，再从后向前移动}
\lstinputlisting[
  language=C++,
  firstline=29,
  lastline=50,
  firstnumber=29,
  numbers=left,
  caption={顺序表插入与删除的核心转移}
]{30_408/10_数据结构/01_线性关系与存储表示/code/ds01_linear_list.hpp}

插入循环必须从 $length$ 向 $index$ 逆向移动；若正向覆盖，尚未复制的旧元素会被破坏。\code{ensure\_capacity} 在移动前执行，保证 $length+1\le capacity$。删除则正向覆盖被删位置，并在移动完成后递减长度。两段代码分别把“不覆盖未迁移数据”和“有效区恰为 $[0,length)$”落成了执行顺序。

扩容本身是一次显式状态迁移：
\lstinputlisting[
  language=C++,
  firstline=76,
  lastline=90,
  firstnumber=76,
  numbers=left,
  caption={容量倍增、复制和所有权切换}
]{30_408/10_数据结构/01_线性关系与存储表示/code/ds01_linear_list.hpp}

旧数据完全复制后才切换 \code{data\_} 与 \code{capacity\_}，因此异常或中途失败不会留下“容量已经变大但数据还没搬完”的半合法状态。一次扩容仍为 $\Theta(n)$；只有分析连续追加序列时，倍增策略才给出摊还 $O(1)$。

\subsection{单链表：局部 \texorpdfstring{$O(1)$}{O(1)} 从“已知前驱”开始}
\lstinputlisting[
  language=C++,
  firstline=142,
  lastline=169,
  firstnumber=142,
  numbers=left,
  caption={已知前驱后的插入与删除}
]{30_408/10_数据结构/01_线性关系与存储表示/code/ds01_linear_list.hpp}

插入先令新节点继承旧后继，再让前驱指向新节点；若交换这两步，新节点会把自己接成错误后继或丢失旧链。删除先保存目标节点，跨过它，再修复 \code{tail\_}、释放节点并递减长度。这里的 $O(1)$ 只从传入合法 \code{predecessor} 开始，不包含寻找前驱。

\subsection{按值删除：Locate 与 Update 在代码中也必须分层}
\lstinputlisting[
  language=C++,
  firstline=171,
  lastline=192,
  firstnumber=171,
  numbers=left,
  caption={按值删除第一次出现的节点}
]{30_408/10_数据结构/01_线性关系与存储表示/code/ds01_linear_list.hpp}

循环负责 Locate，最坏访问全部 $n$ 个节点；找到后才进入头节点分支或复用 \code{erase\_after}。删除唯一节点时必须同时把 \code{head\_} 和 \code{tail\_} 置空，这正是“边界属于合法状态集合”的代码证据。

\subsection{不变量检查与测试不是生产操作的成本}
伴随实现提供 \code{invariant()}：检查空表字段一致性、环、可达节点数、尾节点身份及“尾节点的 \code{next} 为空”。它用于教学验证和测试，不意味着每次生产操作都必须额外遍历全表。

最小断言测试位于 \codepath{code/ds01_linear_list_test.cpp}，覆盖顺序表扩容与非法位置、链表空表、已知前驱插入、首尾删除和删除至空表。测试的关键部分直接验证抽象结果和结构不变量：
\lstinputlisting[
  language=C++,
  firstline=35,
  lastline=57,
  firstnumber=35,
  numbers=left,
  caption={链表边界与不变量断言}
]{30_408/10_数据结构/01_线性关系与存储表示/code/ds01_linear_list_test.cpp}

编译验证命令为：
\begin{lstlisting}[language=bash]
clang++ -std=c++17 -Wall -Wextra -Werror -pedantic \
  code/ds01_linear_list_test.cpp -o /tmp/ds01_test
/tmp/ds01_test
\end{lstlisting}

\begin{methodbox}{做题时的第一判断}
题目给“第 $k$ 个位置/下标”时，先判断表示能否直接映射位置；题目给“某个关键字”时，先把查找成本单列；题目给“节点指针/前驱指针”时，才可以把链式局部修改写成 $O(1)$。不要让操作名称替代输入契约。
\end{methodbox}

\section{空、首、尾与重复值不是补丁}
\begin{center}\small
\begin{tabularx}{\linewidth}{L{2.2cm}YY}
\toprule
边界状态 & 顺序表检查 & 链式表示检查\\
\midrule
空表 & $length=0$，不能访问 $A[0]$ & $head=null$，$tail$ 也必须一致\\
单元素 & 删除后长度变为 0 & 删除后 $head=tail=null$\\
头部 & 插入/删除会移动或更新首位置 & 更新 head，不能丢掉旧链\\
尾部 & 追加可能触发扩容 & 更新 tail，保持 $tail.next=null$\\
重复值 & 明确删第一次、全部还是指定位置 & 不能用“找到一个”替代题目契约\\
越界 & 先判 $0\le k\le length$ & 先判指针/位置可达\\
\bottomrule
\end{tabularx}\end{center}

\begin{warnbox}{最常见的错误推理}
“链表插入一定比数组快”忽略了定位；“数组访问一定 $O(1)$”忽略了按值查找；“长度字段存在就一定正确”忽略了每次分支更新；“删除节点后把指针设空就完成了删除”忽略了前驱、head/tail 和回收责任。
\end{warnbox}

\section{概念边界：把相似说法还原成判据}
\begin{center}\small
\begin{tabularx}{\linewidth}{L{1.8cm}C{0.35cm}L{1.8cm}Y Y}
\toprule
概念 A & $\neq$ & 概念 B & 真正区别与题面信号 & 混淆后果\\
\midrule
逻辑线性表 & $\neq$ & 顺序表 & 前者是关系/操作契约，后者是连续表示 & 把一种实现当成 ADT 定义\\
位置已知 & $\neq$ & 关键字已知 & 位置可直接进入更新，关键字仍需定位 & 漏算查找成本\\
节点已定位 & $\neq$ & 节点可达 & 有指针只说明能沿链走，不说明已拿到前驱 & 把删除写成错误的 $O(1)$\\
长度 & $\neq$ & 容量 & 有效元素数量 vs 已分配槽位数量 & 扩容/越界条件写反\\
访问时间 & $\neq$ & 操作总成本 & 单次寻址/比较 vs 定位、移动、维护的总和 & 复杂度结论失真\\
\bottomrule
\end{tabularx}\end{center}

\section{做题控制协议与压缩复原}
遇到线性表题，按以下顺序写草稿：
\begin{enumerate}
\item 写清逻辑操作的输入契约：位置、关键字、节点还是前驱；
\item 画表示：连续下标或节点/指针，并标出 length、capacity、head、tail；
\item 写操作前不变量；
\item 只执行最小状态更新，再补边界分支；
\item 用一次遍历或边界地址验证逻辑次序、可达性和长度；
\item 最后分别报告定位、局部更新、移动/跳转、额外空间的成本。
\end{enumerate}

\begin{corebox}{一页压缩}
忘掉具体代码后，只需复原五个问题：\kw{对象是什么？输入契约是什么？数据在哪里？操作后必须保持什么？哪一项成本被表示换走了？}
连续表示用地址映射换按位访问；链式表示用指针关系换局部插删。所有“快/慢”判断都必须回到操作输入和不变量。
\end{corebox}

\section{全科坐标与跨册交接}
\begin{corebox}{Map Coordinate：Representation / Operation / Invariant / Cost}
DS01 把 Atlas 的统一语言落到线性关系：输出“逻辑位置—物理地址/指针—可执行操作—合法状态—成本向量”。DS02 复用顺序/链接表示实现端点限制；DS09 复用有序线性存储讨论查找与索引代价；DS-B01 只把链式/数组状态作为 frontier 容器调用。
\end{corebox}
线性表的本册停止点是“表示如何支撑操作”。一旦题目转而问栈/队列为何产生时间秩序、树/图如何展开关系或索引如何维护有序查询，应交给相应 Owner。

\section{来源处理}
本册吸收归档《数据结构 MOC》和《算法时间复杂度：统一分析框架》的对象—表示—操作—成本语言，并将算法笔记中的列表/队列/字典 API 仅作为实现例子，不把 Python 容器行为提升为 408 数据结构定义。扩容摊还、缓存局部性和泛型容器保留为 Extension；栈、队列、串、树、图、Hash 与排序分别由其他 Topic Own。
\end{document}
